\documentclass[11pt]{article}

\usepackage[utf8]{inputenc}
\usepackage[T1]{fontenc}
\usepackage{amsmath,amssymb,amsthm,mathtools}
\usepackage{booktabs}
\usepackage{enumitem}
\usepackage{geometry}
\usepackage{hyperref}
\usepackage{titlesec}
\usepackage[most]{tcolorbox}

\geometry{left=0.9in,right=0.9in,top=0.68in,bottom=0.7in}
\hypersetup{colorlinks=true,linkcolor=blue,urlcolor=blue}
\setlist{itemsep=1pt,topsep=2pt,parsep=0pt,partopsep=0pt}
\setlength{\abovedisplayskip}{6pt plus 2pt minus 2pt}
\setlength{\belowdisplayskip}{6pt plus 2pt minus 2pt}
\setlength{\abovedisplayshortskip}{4pt plus 2pt minus 1pt}
\setlength{\belowdisplayshortskip}{4pt plus 2pt minus 1pt}
\titlespacing*{\section}{0pt}{1.5ex plus .4ex minus .2ex}{.7ex}
\titlespacing*{\subsection}{0pt}{1.2ex plus .3ex minus .2ex}{.5ex}
\titlespacing*{\subsubsection}{0pt}{1ex plus .2ex minus .1ex}{.4ex}

\newcommand{\R}{\mathbb{R}}
\newcommand{\E}{\mathbb{E}}
\newcommand{\trans}{^{\mathsf T}}
\newcommand{\grad}{\nabla}
\newcommand{\norm}[1]{\left\lVert#1\right\rVert}
\newcommand{\argmin}{\operatorname*{arg\,min}}
\newcommand{\proj}{\operatorname{proj}}
\newcommand{\diag}{\operatorname{diag}}

\theoremstyle{definition}
\newtheorem{definition}{Definition}[section]
\theoremstyle{remark}
\newtheorem{remark}[definition]{Remark}

\definecolor{PropositionBlue}{RGB}{31,78,121}
\definecolor{WarningRed}{RGB}{154,27,30}
\newtcbtheorem[number within=section]{keyproposition}{Proposition}{
  enhanced,breakable,
  colback=PropositionBlue!4,
  colframe=PropositionBlue,
  colbacktitle=PropositionBlue,
  coltitle=white,
  fonttitle=\bfseries,
  boxrule=0.9pt,
  arc=1.5mm,
  left=2mm,right=2mm,top=1.2mm,bottom=1.2mm,
  before skip=5pt,after skip=5pt
}{prop}

\newtcolorbox{warningbox}[1]{
  enhanced,breakable,
  title={#1},
  colback=WarningRed!4,
  colframe=WarningRed,
  colbacktitle=WarningRed,
  coltitle=white,
  fonttitle=\bfseries,
  boxrule=0.9pt,
  arc=1.5mm,
  top=1.2mm,bottom=1.2mm,
  before skip=5pt,after skip=5pt
}

\title{Summer Bootcamp -- Session 3 Notes\\Optimization}
\author{Nathan Sauldubois}
\date{August 2026}

\begin{document}
\maketitle

\section*{Purpose of the Session}

Optimization turns a modeling objective into a decision. This session develops
the conditions used to recognize optima, the role of convexity in obtaining
global guarantees, and the algorithms used to solve smooth and constrained
problems. The final sections apply this toolkit to portfolio allocation.

The central logical distinction is the following:
\begin{quote}
derivatives generate candidates; curvature, convexity, feasibility, and
duality determine whether those candidates are optimal.
\end{quote}

\section{Formulating an Optimization Problem}

A general constrained optimization problem has the form
\[
  \min_{x\in\R^d} f(x)
  \quad\text{subject to}\quad
  h(x)=0,\qquad g(x)\leq0.
\]
Here $x$ is the decision variable, $f$ is the objective function, and
\[
  \mathcal F=\{x\in\R^d:h(x)=0,\ g(x)\leq0\}
\]
is the feasible set.

\begin{definition}[Local and global minimizers]
A feasible point $x^\star$ is a local minimizer if it has no better feasible
competitor in some neighborhood. It is a global minimizer if
\[
  f(x^\star)\leq f(x)\qquad\text{for every }x\in\mathcal F.
\]
\end{definition}

Three recurring financial examples are
\[
  \min_\theta \frac12\norm{A\theta-b}_2^2,
\]
\[
  \min_w\frac12w\trans\Sigma w
  \quad\text{s.t.}\quad \mathbf1\trans w=1,
\]
and
\[
  \min_w\frac12w\trans\Sigma w-\gamma\mu\trans w
  \quad\text{s.t.}\quad \mathbf1\trans w=1,\quad w\geq0.
\]

\paragraph{Example.}
Fitting $y_i\approx a+bt_i$ by least squares means choosing
$x=(a,b)\trans$ and solving
\[
  \min_{a,b}\frac12\sum_{i=1}^n(a+bt_i-y_i)^2
  =\min_x\frac12\norm{Ax-y}_2^2,
\]
with $A_i=(1,t_i)$. There is no constraint.

\paragraph{Example.}
If $R$ contains asset returns and $b$ contains benchmark returns, a fully
invested long-only tracking portfolio solves
\[
  \boxed{
  \min_{w\in\R^d}\frac12\norm{Rw-b}_2^2
  \quad\text{s.t.}\quad \mathbf1\trans w=1,\quad w\geq0.}
\]
The decision variable is $w$ and the feasible set is the long-only simplex.

\section{Unconstrained Optimality}

\subsection{First-order condition}

\begin{keyproposition}{First-order necessary condition}{first-order-condition}
Let $f:\R^d\to\R$ be differentiable. If $x^\star$ is an interior local
minimum, then
\[
  \boxed{\grad f(x^\star)=0.}
\]
Such a point is called a critical or stationary point.
\end{keyproposition}

\begin{proof}
For every direction $v\in\R^d$, define
$\phi(t)=f(x^\star+tv)$. Since $t=0$ is a local minimum of $\phi$,
\[
  0=\phi'(0)=\grad f(x^\star)\trans v.
\]
This holds for every $v$, so $\grad f(x^\star)=0$.
\end{proof}

\begin{warningbox}{Necessary does not mean sufficient}
A stationary point can be a minimum, a maximum, or a saddle point. The
first-order equation only generates candidates; classification requires
curvature or another argument.
\end{warningbox}

\subsection{Second-order information}

For a twice differentiable function, Taylor expansion at a critical point gives
\[
  f(x^\star+h)
  =f(x^\star)+\frac12h\trans\grad^2f(x^\star)h+o(\norm{h}^2).
\]

\begin{keyproposition}{Second-order classification}{second-order-classification}
Let $x^\star$ satisfy $\grad f(x^\star)=0$.
\begin{itemize}
  \item If $\grad^2f(x^\star)\succ0$, then $x^\star$ is a strict local minimum.
  \item If $\grad^2f(x^\star)\prec0$, then $x^\star$ is a strict local maximum.
  \item If $\grad^2f(x^\star)$ is indefinite, then $x^\star$ is a saddle point.
  \item A semidefinite Hessian can be inconclusive.
\end{itemize}
\end{keyproposition}

\subsubsection*{Reminder: positivity of symmetric matrices}

For $H=H\trans$, the sign of the quadratic form $z\trans Hz$ determines its
definiteness:
\[
\boxed{
\begin{aligned}
H\succ0&\Longleftrightarrow z\trans Hz>0 &&\text{for all }z\neq0,\\
H\succeq0&\Longleftrightarrow z\trans Hz\geq0 &&\text{for all }z,\\
H\prec0&\Longleftrightarrow z\trans Hz<0 &&\text{for all }z\neq0.
\end{aligned}}
\]
If the quadratic form takes both signs, $H$ is indefinite. For symmetric
matrices, these conditions are equivalent to the corresponding sign conditions
on every eigenvalue.

For a $2\times2$ symmetric matrix
\[
H=\begin{pmatrix}a&b\\b&c\end{pmatrix},
\]
the easy determinant tests are
\[
\boxed{
\begin{array}{rcl}
H\succ0&\Longleftrightarrow&a>0\text{ and }ac-b^2>0,\\
H\prec0&\Longleftrightarrow&a<0\text{ and }ac-b^2>0,\\
H\text{ indefinite}&\Longleftrightarrow&ac-b^2<0.
\end{array}}
\]
In any dimension, Sylvester's criterion states that $H\succ0$ exactly when
all leading principal minors are positive.

\paragraph{Example.}
For $f(x)=x^3-3x^2+2$, the stationary points are $0$ and $2$. Since
$f''(0)=-6$ and $f''(2)=6$, they are respectively a local maximum and a local
minimum. For
\[
q(x,y)=x^2+xy+y^2-4x-2y,
\]
the stationary point is $(2,0)$ and
\[
H_q=\begin{pmatrix}2&1\\1&2\end{pmatrix}\succ0,
\]
so it is the unique global minimum. Finally, $x^2-y^2$ has an indefinite
Hessian at the origin and therefore a saddle point there.

\paragraph{Example.}
For
\[
g_1(x,y)=3x^2+2xy+2y^2,
\qquad g_2(x,y)=-x^2-xy-y^2,
\]
we have
\[
H_{g_1}=\begin{pmatrix}6&2\\2&4\end{pmatrix}\succ0,
\qquad
H_{g_2}=\begin{pmatrix}-2&-1\\-1&-2\end{pmatrix}\prec0.
\]
Thus the origin is a strict minimum of $g_1$ and a strict maximum of $g_2$.

The functions $x\mapsto x^4$ and $x\mapsto-x^4$ both have zero first and
second derivatives at the origin, yet the origin is respectively a strict
minimum and a strict maximum. Higher-order information is then necessary.

\subsection{Quadratic objectives and least squares}

Consider
\[
  f(x)=\frac12x\trans Qx+c\trans x+r,
  \qquad Q=Q\trans.
\]
Then
\[
  \grad f(x)=Qx+c,
  \qquad \grad^2f(x)=Q.
\]

\begin{keyproposition}{Minimizer of a positive-definite quadratic}{quadratic-minimizer}
If $Q\succ0$, the function has the unique global minimizer
\[
  \boxed{x^\star=-Q^{-1}c.}
\]
\end{keyproposition}

For least squares,
\[
  f(x)=\frac12\norm{Ax-b}_2^2,
  \qquad
  \grad f(x)=A\trans(Ax-b).
\]
The criticality condition produces the normal equations
\[
  \boxed{A\trans A x=A\trans b.}
\]
If $A$ has full column rank, then $A\trans A\succ0$ and the solution is unique.

\section{Convexity and Global Optimality}

\begin{definition}[Convex set]
A set $C\subset\R^d$ is convex if
\[
  x,y\in C,\ t\in[0,1]
  \quad\Longrightarrow\quad
  tx+(1-t)y\in C.
\]
\end{definition}

\begin{definition}[Convex function]
A function $f:C\to\R$ is convex if
\[
  f(tx+(1-t)y)\leq tf(x)+(1-t)f(y)
\]
for all $x,y\in C$ and $t\in[0,1]$.
\end{definition}

\begin{keyproposition}{First-order characterization of convexity}{convexity-first-order}
Let $f$ be differentiable on an open convex set. Then $f$ is convex if and only if
\[
  \boxed{f(y)\geq f(x)+\grad f(x)\trans(y-x)
  \qquad\text{for all }x,y.}
\]
Thus the tangent affine function is a global lower bound.
\end{keyproposition}

\begin{keyproposition}{Criticality is sufficient under convexity}{convex-criticality}
If $f$ is differentiable and convex and
\[
  \grad f(x^\star)=0,
\]
then $x^\star$ is a global minimizer.
\end{keyproposition}

Indeed, the first-order inequality gives $f(y)\geq f(x^\star)$ for every $y$.

\begin{keyproposition}{Second-order characterization of convexity}{convexity-second-order}
Let $f$ be twice differentiable on an open convex set. Then
\[
  \boxed{f\text{ is convex}
  \quad\Longleftrightarrow\quad
  \grad^2f(x)\succeq0\quad\text{for every }x.}
\]
\end{keyproposition}

\paragraph{Example.}
The functions
\[
e^{a\trans x},\qquad \log(1+e^x),\qquad x^2-xy+y^2
\]
are convex because their Hessians are respectively
\[
e^{a\trans x}aa\trans\succeq0,
\qquad
\frac{e^x}{(1+e^x)^2}>0,
\qquad
\begin{pmatrix}2&-1\\-1&2\end{pmatrix}\succ0.
\]

\paragraph{Example.}
For $p(x,y)=x^4+e^y$,
\[
\grad^2p(x,y)=\begin{pmatrix}12x^2&0\\0&e^y\end{pmatrix}\succeq0,
\]
so $p$ is convex. For $q(x,y)=x^2-y^2$, the Hessian
$\diag(2,-2)$ is indefinite, so $q$ is not convex.

\begin{definition}[Strong convexity and smoothness]
A twice differentiable function is $\mu$-strongly convex and $L$-smooth when
\[
  \mu I\preceq\grad^2f(x)\preceq LI
\]
throughout its domain. The ratio $\kappa=L/\mu$ is the condition number.
\end{definition}

Large $\kappa$ corresponds to very different curvatures in different
directions. This geometry typically slows first-order algorithms.

For example, the Hessian of $(x^2+100y^2)/2$ has eigenvalues $1$ and $100$,
so $(\mu,L,\kappa)=(1,100,100)$. As a second calculation, the Hessian of
$2x^2+32y^2$ is $\diag(4,64)$, hence
\[
\boxed{\mu=4,\qquad L=64,\qquad\kappa=16.}
\]

\section{Algorithms for Smooth Optimization}

\subsection{Gradient descent}

Gradient descent uses the iteration
\[
  \boxed{x_{k+1}=x_k-\alpha_k\grad f(x_k).}
\]
The negative gradient is the direction of steepest local decrease in the
Euclidean geometry.

\begin{keyproposition}{Descent lemma}{descent-lemma}
If $\grad f$ is $L$-Lipschitz, then
\[
  f(x-\alpha\grad f(x))
  \leq f(x)-\alpha\left(1-\frac{L\alpha}{2}\right)
  \norm{\grad f(x)}_2^2.
\]
Consequently, every fixed step satisfying
\[
  \boxed{0<\alpha<\frac{2}{L}}
\]
decreases the objective away from stationary points.
\end{keyproposition}

A standard conservative choice is $\alpha=1/L$. When $L$ is unknown,
backtracking line search reduces the step until sufficient decrease is observed.

\paragraph{Example.}
For $f(x)=\frac12(x-3)^2$, $x_0=0$, and $\alpha=0.4$,
\[
x_{k+1}=0.6x_k+1.2,
\qquad (x_1,x_2,x_3)=(1.2,1.92,2.352).
\]
For the companion problem $g(x)=\frac12(x+2)^2$, $x_0=2$, and
$\alpha=0.25$,
\[
x_{k+1}=0.75x_k-0.5,
\qquad
\boxed{(x_1,x_2,x_3)=(1,0.25,-0.3125).}
\]

\begin{keyproposition}{Linear convergence under strong convexity}{gradient-linear-rate}
If $f$ is $\mu$-strongly convex and $L$-smooth, gradient descent with
$\alpha=1/L$ satisfies
\[
  \boxed{
  f(x_k)-f(x^\star)
  \leq
  \left(1-\frac{\mu}{L}\right)^k
  \bigl(f(x_0)-f(x^\star)\bigr).}
\]
\end{keyproposition}

\subsection{Newton's method}

The second-order approximation
\[
  f(x+p)\approx f(x)+\grad f(x)\trans p
  +\frac12p\trans\grad^2f(x)p
\]
is minimized by solving
\[
  \grad^2f(x_k)p_k=-\grad f(x_k),
  \qquad x_{k+1}=x_k+p_k.
\]

\begin{warningbox}{Implementation rule}
Solve the linear system for $p_k$; do not explicitly form the inverse Hessian.
Globalization may require damping or a line search, especially when the
Hessian is indefinite or the initial point is far from the solution.
\end{warningbox}

\paragraph{Example.}
For $f(x)=e^x-3x$, Newton's update is
\[
x_{k+1}=x_k-1+3e^{-x_k}.
\]
Starting from $x_0=1$ gives $x_1\simeq1.10364$ and
$x_2\simeq1.09862$, close to $\log3$. For the companion problem
$g(x)=e^x-2x$ from $x_0=0.5$,
\[
x_{k+1}=x_k-1+2e^{-x_k},
\qquad
\boxed{x_1\simeq0.71306,\quad x_2\simeq0.69334,}
\]
close to $\log2$.

\subsection{Projected gradient descent}

We now solve the convex constrained problem
\[
  \min_{x\in C}f(x),
\]
where $f$ is differentiable and convex and $C$ is nonempty, closed, and convex.

\begin{keyproposition}{Euclidean projection onto a convex set}{convex-projection}
For every $z\in\R^d$, the projection
\[
  \boxed{\proj_C(z)=\argmin_{u\in C}\norm{u-z}_2}
\]
exists and is unique.
\end{keyproposition}

Projected gradient descent alternates an unconstrained step and a projection:
\[
  \boxed{x_{k+1}=\proj_C\bigl(x_k-\alpha_k\grad f(x_k)\bigr).}
\]

\section{Equality Constraints and Lagrange Multipliers}

Consider
\[
  \min_x f(x)\quad\text{s.t.}\quad h(x)=0.
\]
At a regular feasible point, first-order feasible directions satisfy
$\grad h(x)\trans v=0$. At an optimum, $\grad f(x^\star)$ must also be
orthogonal to all feasible directions.

\begin{keyproposition}{Lagrange first-order conditions}{lagrange-conditions}
Define
\[
  \mathcal L(x,\lambda)=f(x)+\lambda\trans h(x).
\]
Under a regularity condition, a constrained local optimum satisfies
\[
  \boxed{
  \grad_x\mathcal L(x^\star,\lambda^\star)=0,
  \qquad h(x^\star)=0.}
\]
\end{keyproposition}

These equations are necessary in general, not sufficient. Convexity or a
second-order constrained argument is required to classify candidates.

\paragraph{Example.}
For
\[
\min_{x,y}x^2+y^2\quad\text{s.t.}\quad x+y=1,
\]
the equations $2x+\lambda=0$, $2y+\lambda=0$, and $x+y=1$ give
\[
\boxed{x^\star=y^\star=\frac12,\qquad\lambda^\star=-1.}
\]

\paragraph{Example.}
For
\[
\min_{x,y}x^2+y^2\quad\text{s.t.}\quad2x-y=3,
\]
the Lagrangian is $x^2+y^2+\lambda(2x-y-3)$. The first-order equations give
$x=-\lambda$, $y=\lambda/2$, and therefore
\[
\boxed{\lambda^\star=-\frac65,\qquad
(x^\star,y^\star)=\left(\frac65,-\frac35\right).}
\]

For several equality constraints $h=(h_1,\ldots,h_m)$, a standard regularity
condition is LICQ.

\begin{keyproposition}{Linear independence constraint qualification}{licq}
LICQ holds at $x^\star$ when
\[
  \boxed{\grad h_1(x^\star),\ldots,\grad h_m(x^\star)
  \text{ are linearly independent}.}
\]
\end{keyproposition}

\paragraph{Why regularity matters.}
For $\min_x x$ subject to $x^2=0$, the only feasible point is $0$, but
$\partial_x(x+\lambda x^2)|_{x=0}=1$ for every $\lambda$. Likewise, for
\[
\min_{x,y}x+y\quad\text{s.t.}\quad x^2+y^2=0,
\]
the only feasible point is $(0,0)$ while
$\grad f(0,0)=(1,1)\trans$ and $\grad h(0,0)=0$. In both cases the constraint
gradient vanishes, LICQ fails, and no multiplier can satisfy stationarity.

\subsection{Duality}

The multiplier $\lambda_i$ can be interpreted as the marginal price of
violating constraint $i$. The dual function is
\[
  q(\lambda)=\inf_x\mathcal L(x,\lambda).
\]
It is concave because it is the pointwise infimum of affine functions of
$\lambda$.

\begin{keyproposition}{Weak duality}{weak-duality}
Let
\[
  p^\star=\inf_{h(x)=0}f(x),
  \qquad d^\star=\sup_\lambda q(\lambda).
\]
Then, without requiring convexity,
\[
  \boxed{d^\star\leq p^\star.}
\]
For every feasible $x$ and every $\lambda$,
\[
  \boxed{f(x)-q(\lambda)\geq0.}
\]
The latter quantity is the primal--dual gap.
\end{keyproposition}

A small primal--dual gap gives a computable certificate of near-optimality.

\begin{keyproposition}{Strong duality and saddle points}{strong-duality}
Strong duality means
\[
  \boxed{d^\star=p^\star.}
\]
For convex problems, affine equality constraints and a suitable constraint
qualification commonly ensure this equality. If optimal solutions exist, then
\[
  \boxed{
  \min_x\max_\lambda\mathcal L(x,\lambda)
  =\max_\lambda\min_x\mathcal L(x,\lambda).}
\]
The optimal pair is a saddle point:
\[
  \mathcal L(x^\star,\lambda)
  \leq\mathcal L(x^\star,\lambda^\star)
  \leq\mathcal L(x,\lambda^\star).
\]
\end{keyproposition}

\subsection{Primal--dual gradient method}

For a differentiable Lagrangian, the direct descent--ascent iteration is
\[
  \boxed{
  \begin{aligned}
  x_{k+1}&=x_k-\alpha_k\grad_x\mathcal L(x_k,\lambda_k),\\
  \lambda_{k+1}&=\lambda_k+\beta_k\grad_\lambda
  \mathcal L(x_k,\lambda_k).
  \end{aligned}}
\]
Since $\grad_\lambda\mathcal L(x,\lambda)=h(x)$,
\[
  \lambda_{k+1}=\lambda_k+\beta_k h(x_k).
\]
The primal step decreases the Lagrangian with respect to the decision, while
the dual step raises the price of constraint violations.

\begin{warningbox}{Possible oscillations}
Simultaneous gradient descent--ascent can rotate around a saddle point rather
than converge. Damping, extragradient methods, and augmented Lagrangians can
improve stability.
\end{warningbox}

\section{Inequality Constraints and KKT Conditions}

Consider
\[
  \min_x f(x)
  \quad\text{s.t.}\quad h(x)=0,\qquad g(x)\leq0,
\]
with
\[
  \mathcal L(x,\nu,\lambda)
  =f(x)+\nu\trans h(x)+\lambda\trans g(x).
\]

\begin{keyproposition}{Karush--Kuhn--Tucker necessary conditions}{kkt-necessary}
Under a suitable constraint qualification, a local optimum admits multipliers
such that
\[
  \boxed{
  \begin{aligned}
  \grad_x\mathcal L(x^\star,\nu^\star,\lambda^\star)&=0
    &&\text{stationarity},\\
  h(x^\star)=0,\quad g(x^\star)&\leq0
    &&\text{primal feasibility},\\
  \lambda^\star&\geq0
    &&\text{dual feasibility},\\
  \lambda_i^\star g_i(x^\star)&=0
    &&\text{complementary slackness}.
  \end{aligned}}
\]
\end{keyproposition}

Complementary slackness separates active and inactive constraints:
\begin{itemize}
  \item if $g_i(x^\star)<0$, then $\lambda_i^\star=0$;
  \item if $\lambda_i^\star>0$, then $g_i(x^\star)=0$;
  \item an active constraint may still have a zero multiplier.
\end{itemize}

\paragraph{Example.}
For
\[
\min_{x,y}(x-1)^2+(y-2)^2
\quad\text{s.t.}\quad x+y\leq1,
\]
the unconstrained minimizer is infeasible, so the constraint binds. The KKT
equations give
\[
\boxed{\lambda^\star=2,\qquad(x^\star,y^\star)=(0,1).}
\]

\paragraph{Example.}
For
\[
\min_{x,y}(x-2)^2+(y+1)^2
\quad\text{s.t.}\quad x-y\leq1,
\]
the unconstrained point $(2,-1)$ is infeasible. With multiplier $\lambda$,
stationarity gives $x=2-\lambda/2$ and $y=-1+\lambda/2$. Binding feasibility
then gives
\[
\boxed{\lambda^\star=2,\qquad(x^\star,y^\star)=(1,0).}
\]

\begin{keyproposition}{KKT sufficiency for convex problems}{kkt-sufficient}
Assume that $f$ and every $g_i$ are convex, every equality constraint is
affine, and a constraint qualification holds. Then every point satisfying the
KKT conditions is a global optimum.
\end{keyproposition}

\section{Portfolio Optimization}

Let $R\in\R^d$ be the vector of asset returns, with
\[
  \mu=\E[R],
  \qquad \Sigma=\operatorname{Cov}(R).
\]
For portfolio weights $w$,
\[
  \E[w\trans R]=\mu\trans w,
  \qquad
  \operatorname{Var}(w\trans R)=w\trans\Sigma w.
\]
Since covariance matrices are positive semidefinite, portfolio variance is a
convex quadratic function.

Throughout this section,
\[
\boxed{\mathbf1=(1,\ldots,1)\trans\in\R^d}
\]
denotes the vector of ones, so $\mathbf1\trans w=1$ is the full-investment
constraint.

\subsection{Minimum variance}

\begin{keyproposition}{Fully invested minimum-variance portfolio}{minimum-variance-portfolio}
Assume $\Sigma\succ0$. The problem
\[
  \min_w\frac12w\trans\Sigma w
  \quad\text{s.t.}\quad \mathbf1\trans w=1
\]
has the unique solution
\[
  \boxed{
  w^\star=
  \frac{\Sigma^{-1}\mathbf1}
  {\mathbf1\trans\Sigma^{-1}\mathbf1}.}
\]
\end{keyproposition}

Indeed, stationarity gives $\Sigma w+\lambda\mathbf1=0$, and the budget
constraint determines the multiplier.

\paragraph{Example.}
For $\Sigma=\diag(0.04,0.09,0.16)$, inverse-variance weighting gives
\[
w^\star=\begin{pmatrix}36/61\\16/61\\9/61\end{pmatrix}.
\]
For the companion covariance $\Sigma=\diag(0.01,0.04,0.09)$, the inverse
variances are $100$, $25$, and $100/9$, hence
\[
\boxed{w^\star=\begin{pmatrix}36/49\\9/49\\4/49\end{pmatrix}.}
\]
In both cases the least volatile asset receives the largest weight.

\subsection{Mean--variance trade-off}

A common formulation is
\[
  \boxed{
  \min_w\frac12w\trans\Sigma w-\gamma\mu\trans w
  \quad\text{s.t.}\quad\mathbf1\trans w=1.}
\]
The parameter $\gamma>0$ controls the reward assigned to expected return.
The first-order system is
\[
  \Sigma w-\gamma\mu+\lambda\mathbf1=0,
  \qquad \mathbf1\trans w=1.
\]

Fixing a target return $m$ gives
\[
  \min_w\frac12w\trans\Sigma w
  \quad\text{s.t.}\quad
  \mathbf1\trans w=1,
  \quad \mu\trans w=m.
\]
Let
\[
  C=\begin{pmatrix}\mathbf1\trans\\\mu\trans\end{pmatrix},
  \qquad d=\begin{pmatrix}1\\m\end{pmatrix}.
\]

\begin{keyproposition}{Equality-constrained efficient portfolio}{efficient-portfolio}
If the required inverses exist, the target-return portfolio is
\[
  \boxed{
  w^\star=\Sigma^{-1}C\trans
  (C\Sigma^{-1}C\trans)^{-1}d.}
\]
Varying $m$ traces the mean--variance frontier.
\end{keyproposition}

\subsection{Long-only and operational constraints}

With $w\geq0$, the Lagrangian includes multipliers $\eta\geq0$ and KKT gives
\[
  \Sigma w-\gamma\mu+\lambda\mathbf1-\eta=0,
  \qquad \eta_iw_i=0.
\]
Because the problem is convex, KKT conditions are sufficient under the usual
constraint qualification.

Real portfolio models may also impose position bounds, sector limits,
turnover constraints, transaction costs, tracking-error limits, and stress
constraints. Quadratic risk combined with linear constraints remains a convex
quadratic program and can be solved reliably at realistic scale.

\subsubsection*{Constrained rebalancing as a convex QCQP}

Consider
\[
\min_w\frac12w\trans\Sigma w-\gamma\mu\trans w
\quad\text{s.t.}\quad
\mathbf1\trans w=1,\quad w\geq0,\quad
\norm{w-w^{\mathrm{old}}}_2^2\leq\tau^2.
\]
Using multipliers $\lambda\in\R$, $\eta\geq0$, and $\rho\geq0$, stationarity is
\[
\boxed{
\Sigma w-\gamma\mu+\lambda\mathbf1-\eta
+2\rho(w-w^{\mathrm{old}})=0.}
\]
Complementary slackness requires
\[
\boxed{
\eta_iw_i=0\quad\forall i,
\qquad
\rho\bigl(\norm{w-w^{\mathrm{old}}}_2^2-\tau^2\bigr)=0.}
\]
Together with primal and dual feasibility, these are the KKT conditions. When
$\Sigma\succeq0$, the problem is a convex quadratically constrained quadratic
program. Under the usual qualification, KKT is sufficient and a convex QCQP
solver returns the global optimum.

\section{Summary}

The workflow is to identify the decision variable, objective, and feasible set;
check differentiability and convexity; derive and classify the first-order
conditions; state the sign convention and all KKT conditions when constraints
are present; choose an algorithm adapted to the structure; verify feasibility,
stationarity, and numerical convergence; and interpret the optimizer in the
original application.

\end{document}
